\section{Experimentelle Tragfähigkeitsermittlung und Statistik}
\label{sec:experiment_stand}

Die Werkstoffkennwerte der Berechnung nach VDI 2736 Blatt 2, darunter die Zeitschwellfestigkeit $\sigma_{\mathrm{FlimN}}$, werden nach VDI 2736 Blatt 4 durch Versuche an realen Zahnrädern bestimmt \cite[Abschnitt~3, S.~5]{VDI2736-4}. Dazu laufen die Räder auf Prüfstandsgetrieben bei festgehaltener Werkstofftemperatur mit konstanter Drehzahl und konstantem Drehmoment. Gezählt wird die Zahl der Lastwechsel bis zum Versagen des Bauteils \cite[S.~6]{VDI2736-4}. Vor den Messungen werden Temperatur-Drehmoment-Paare festgelegt, wobei für mindestens eine Temperatur und mindestens vier Drehmomente je ein Versuchsblock mit mindestens drei Wiederholungen vorzusehen ist \cite[Abschnitt~3.1, S.~6]{VDI2736-4}. Als Versagen gilt der Bruch des Zahns oder das Auftreten einer bestimmten Menge an Grübchen. Für welche Schadensart ein Lauf zählt, lässt sich mitunter erst nach seinem Ende entscheiden \cite[S.~6]{VDI2736-4}. Aus dem Drehmoment folgt die Zeitschwellfestigkeit analog zu Gleichung (10) von Blatt 2, und die gemessenen Lastwechselzahlen werden auf $10$\,\% Schadenswahrscheinlichkeit umgerechnet \cite[Abschnitt~3.3, S.~7]{VDI2736-4}. Das Tripel aus Festigkeit, Lastwechselzahl bei $10$\,\% und Prüftemperatur bildet einen Punkt der temperaturabhängigen Wöhlerlinie \cite[S.~7]{VDI2736-4}. Für den ersten Versuch empfiehlt die Richtlinie ein hohes Startmoment mit $10^{5}$ bis $10^{6}$ Lastwechseln, weil zu kleine Momente Durchläufer erzeugen und die Versuchsdauer verlängern \cite[Abschnitt~8.4, S.~20]{VDI2736-4}.

Beim Verspannungsprinzip nach Blatt 4 wird die Prüfverzahnung mit einem Übertragungsgetriebe statisch verspannt. Der Prüfstand soll drehweich sein und das Verspannmoment im Betrieb messen, weil Verschleiß oder Verformung der Kunststoffverzahnung das Moment abfallen lassen \cite[Abschnitt~5.2, S.~11]{VDI2736-4}. Hasl betreibt die Paarungen in einstufigen Wöhlerversuchen bei konstanter Zahntemperatur bis zum Ausfall durch Loss of Drive, einen augenblicklichen Abfall des gemessenen Lastmoments \cite[Abschnitt~7.2, S.~91]{Has18}. Weil das Verspannmoment durch Setzeffekte über der Laufzeit geringfügig sinkt, wertet er das vor dem Ausfall gemessene Moment aus \cite[S.~91]{Has18}. Fiederling beobachtet bei Zahnfußbrüchen ein schleichendes Risswachstum ohne eindeutigen Ausfallzeitpunkt. Neben dem vollständigen Bruch definiert sie deshalb einen Abfall des Drehmoments von mehr als $1$\,Nm über 100\,000 Lastwechsel als Ausfallkriterium \cite[Abschnitt~6.1, S.~65]{Fie23}. Die Temperatur wird berührungslos durch Thermografie oder berührend mit Thermoelementen gemessen, wobei der temperaturabhängige Emissionskoeffizient des Werkstoffs in Referenzmessungen zu bestimmen ist \cite[Abschnitt~8.2, S.~16]{VDI2736-4}. Bei Ölschmierung genügt näherungsweise die Öltemperatur als Maß für die Zahnfußtemperatur \cite[Abschnitt~8.2.2, S.~16]{VDI2736-4}. Öl erschwert die optische Messung, und die Zuverlässigkeit des Zahnflankentemperaturkennwerts ist für Stirnräder nicht nachgewiesen \cite[S.~16 und 18]{VDI2736-4}.

% TK erledigt (alle Stellen geprüft, keine männlichen Pronomen mehr): Fiederling ist weiblich, bitte überall darauf achten.

Der Pulsatorversuch belastet einen oder zwei Zähne mit einer schwellenden Kraft, braucht kein Gegenrad und erlaubt mehrere Versuche je Rad \cite[Abschnitt~2.2.2, S.~12]{Fie23}. Für Kunststoffräder verlangt Blatt 4 eine korrekte Temperierung, weil die im Zahnfuß durch Verformungsenergie entstehende Wärme die Ergebnisse stark beeinflusst, sowie einen fest definierten Kraftangriffspunkt \cite[Abschnitt~5.3, S.~12]{VDI2736-4}. Die Übertragbarkeit ist durch vergleichende Laufversuche zu prüfen \cite[S.~12]{VDI2736-4}. Hasl nennt als Störgrößen den gegenüber dem Laufbetrieb veränderten zeitlichen Beanspruchungsverlauf, die Haltekraft mit dem Kriechen des eingespannten Zahns, die dehnratenabhängige Steifigkeit und die innere Erwärmung \cite[Abschnitt~5.1, S.~61 bis 62]{Has18}. Der Aufwand zu ihrer Beherrschung liegt über dem des Laufversuchs, der zudem dem Anwendungsfall entspricht \cite[S.~62]{Has18}. Fiederling nennt als Unterschiede das Versagen des schwächeren von zwei Zähnen gegenüber dem schwächsten von 52, das Spannungsverhältnis $R > 0$ gegenüber $R \le 0$ und die feste gegenüber der wandernden Berührlinie \cite[Abschnitt~2.2.3, S.~14 bis 15]{Fie23}. Umrechnungsfaktoren zwischen beiden Methoden liegen bislang nur für Stahlräder vor \cite[S.~15]{Fie23}. Für die Verzahnung kst-E aus POM bei $80$\,\textdegree C liegen die Wöhlerlinien des Pulsators links von denen des Laufversuchs und verlaufen steiler \cite[Abschnitt~8.1, S.~105 bis 106]{Fie23}. Im Hauptlastbereich von $12$\,Nm ergibt sich ein überschlägiger Korrelationsfaktor der Lebensdauern von etwa $2{,}6$ \cite[S.~106]{Fie23}.

Die Zeitschwellfestigkeit nach Blatt 2 hängt von der Zahnfußtemperatur und der Lastwechselzahl $N_{\mathrm{L}}$ ab \cite[Abschnitt~3.2.2, S.~14]{VDI2736-2}. Bild 3 und Tabelle 5 geben sie für PA66 und POM von $10^{5}$ bis $10^{8}$ Lastwechseln an, und bei Dauerbetrieb ist mit $N_{\mathrm{L}} = 10^{8}$ zu rechnen \cite[S.~14 bis 15]{VDI2736-2}. Für POM bei $80$\,\textdegree C fällt der Wert von $50$\,N/mm$^{2}$ bei $10^{5}$ auf $20$\,N/mm$^{2}$ bei $10^{8}$ Lastwechseln. Ein Abknicken in eine Dauerfestigkeit sieht die Richtlinie nicht vor \cite[Tabelle~5, S.~15]{VDI2736-2}. Hasl übernimmt die Darstellung der Richtlinie mit linearer Spannungsachse über der logarithmischen Lastspielzahl und zeigt, dass die doppeltlogarithmische Darstellung im Beispiel keine gravierenden Unterschiede ergibt \cite[S.~107 und 181]{Has18}. Ob bei vernachlässigbar kleiner plastischer Verformung eine technisch relevante Dauerfestigkeit thermoplastischer Kunststoffe existiert, lässt er offen \cite[S.~176]{Has18}. Fiederling findet für POM der Verzahnung kst-E eine Wöhlerlinie, die sich bis etwa $40 \cdot 10^{6}$ Lastwechsel bei $8$\,Nm mit unveränderter Steigung verlängern lässt, und keine Dauerfestigkeit \cite[Abschnitt~6.3.1, S.~82 bis 83]{Fie23}. Ihre Ergebnisse stützen die dort zitierte These einer linearen Fortführung der Zeitfestigkeitsgeraden bis mindestens $10^{10}$ Lastwechsel \cite[S.~83]{Fie23}. Stommel et al. kennen für Kunststoffe keine Dauerfestigkeitsgrenze und wählen als Ersatz die Spannungsamplitude bei der Grenzlastspielzahl $N_{\mathrm{G}} = 10^{7}$, eine willkürliche Auslegungsgrenze, jenseits derer mit weiterem Abfall zu rechnen ist \cite[S.~328]{Sto18}. Den Zeitfestigkeitsbereich beschreiben sie mit der Basquin-Gleichung, deren doppeltlogarithmische Gerade mit dem Neigungsexponenten $k$ auf die Grenzspannungsamplitude bezogen ist \cite[S.~206 bis 207]{Sto18}. DIN 50100 ersetzt den Begriff der Dauerfestigkeit durch die Langzeitfestigkeit, die sich auf die Grenzschwingspielzahl bezieht, ab der ein Versuch als Durchläufer zählt \cite[Abschnitt~5, S.~16]{DIN50100}. Ohne experimentell belegte Langzeitfestigkeit wird die Zeitfestigkeitsgerade nur zwischen dem höchsten und dem niedrigsten belegten Lasthorizont dargestellt \cite[S.~16]{DIN50100}.

DIN 50100 regelt Durchführung und Auswertung von Schwingfestigkeitsversuchen mit konstanter Lastamplitude an metallischen Werkstoffproben und Bauteilen zwischen $10^{4}$ und $10^{7}$ Schwingspielen \cite[Abschnitt~1, S.~9]{DIN50100}. Sie gibt die Wöhlerlinie für $50$\,\% Ausfallwahrscheinlichkeit sowie Linien für $10$\,\% und $90$\,\% an und verwendet ausschließlich die logarithmische Normalverteilung \cite[S.~9]{DIN50100}. Im Bereich üblicher Stichproben von 10 bis 25 Proben unterscheiden sich die gebräuchlichen Verteilungsfunktionen nur gering. Erst bei sehr kleinen Ausfallwahrscheinlichkeiten laufen sie auseinander \cite[Abschnitt~5, S.~18]{DIN50100}. Während der Mittelwert schon mit kleinen Stichproben gut geschätzt wird, verlangt die Standardabweichung sehr große Stichprobenumfänge \cite[S.~19 und 54]{DIN50100}. Für die Zeitfestigkeitsgerade sieht die Norm das Perlenschnur- oder das Horizontenverfahren vor. Beim Perlenschnurverfahren beginnt man auf einem beliebigen Lasthorizont und tastet sich mit den folgenden Versuchen an die Übergangsbereiche heran, beim Horizontenverfahren werden vorab zwei weit auseinanderliegende Horizonte festgelegt \cite[Abschnitt~9.2.1, S.~26]{DIN50100}. Die Versuche des Perlenschnurverfahrens werden gemeinsam durch eine Regression nach der Methode der kleinsten Fehlerquadrate in Richtung der Schwingspielzahl ausgeglichen, wobei mehrere Versuche auf einem Horizont zulässig sind \cite[Abschnitt~9.2.2, S.~28]{DIN50100}. Liegen Ergebnisse auf mehr als zwei Horizonten vor, werden sie auch im Horizontenverfahren auf diese Weise ausgewertet \cite[Abschnitt~10.2.3, S.~57]{DIN50100}. \figref{fig:din50100_perlenschnur} zeigt das Schema.

\begin{figure}[htbp]
    \centering
    \includegraphics[width=0.8\textwidth]{kap2/quelle_din50100_perlenschnur}
    \caption{Schema des Perlenschnurverfahrens im Wöhlerdiagramm, nach \cite[Bild~8, S.~29]{DIN50100}.}
    \label{fig:din50100_perlenschnur}
\end{figure}

Die Zeitfestigkeitsgerade folgt dem Ansatz von Basquin, der durch Logarithmieren in eine Geradengleichung übergeht \cite[Abschnitt~10.2.2, S.~55]{DIN50100}:
\begin{equation}
    \lg N_{\mathrm{L}} = \lg C - k \,\lg L_{\mathrm{a}} ,
    \label{eq:zeitfestigkeitsgerade}
\end{equation}
mit dem Lageparameter $C$, der Neigung $k$ und der Lastamplitude $L_{\mathrm{a}}$. Beide Konstanten folgen aus der linearen Regression mit der Schwingspielzahl als abhängiger Größe. Der Berechnung der Standardabweichung liegt die Annahme zugrunde, dass sie auf allen Lasthorizonten gleich ist, weshalb die Ergebnisse parallel zur Ausgleichsgeraden auf einen fiktiven Horizont verschoben werden \cite[S.~56]{DIN50100}:
\begin{equation}
    \begin{aligned}
        N_{i,\mathrm{fiktiv}} & = N_{i} \left(\frac{L_{\mathrm{a,fiktiv}}}{L_{\mathrm{a},i}}\right)^{-k} ,                                               \\
        \tilde{s}_{\lg N}     & = \sqrt{\frac{1}{n-2} \sum_{i=1}^{n} \left(\lg N_{i,\mathrm{fiktiv}} - \lg N_{50\%,\mathrm{fiktiv}}\right)^{2}} , \qquad
        \tilde{s}_{\lg N,\mathrm{korr}} = \tilde{s}_{\lg N} \, \frac{n - 1{,}74}{n - 2} .
    \end{aligned}
    \label{eq:perlenschnur_streuung}
\end{equation}
Die zwei verlorenen Freiheitsgrade entfallen auf die geschätzte Lage und die geschätzte Neigung. Weil kleine Stichproben die Standardabweichung im Mittel unterschätzen, wird sie mit dem Faktor der dritten Gleichung korrigiert \cite[S.~56 bis 57]{DIN50100}. Für eine Ausfallwahrscheinlichkeit $P_{\mathrm{A}}$ folgt die Schwingspielzahl aus dem Mittelwert, der korrigierten Standardabweichung und dem Quantil $u$ der Standardnormalverteilung \cite[Abschnitt~10.5, S.~63]{DIN50100}:
\begin{equation}
    N_{P_{\mathrm{A}}} = 10^{\lg N_{50\%} + u \, s_{\lg N,\mathrm{korr}}} .
    \label{eq:quantil}
\end{equation}
Tabelle 11 der Norm nennt für $10$\,\% das Quantil $u = -1{,}282$ und für $90$\,\% den Wert $u = 1{,}282$ \cite[Tabelle~11, S.~63]{DIN50100}. Im Perlenschnurverfahren gilt die Standardabweichung entlang der Zeitfestigkeitsgeraden als konstant, die Linien verschiedener Ausfallwahrscheinlichkeiten verlaufen parallel \cite[S.~63]{DIN50100}. Übliche Standardabweichungen $s_{\lg N}$ liegen zwischen $0{,}10$ und $0{,}30$. Ohne Erfahrungswerte wird $0{,}20$ angenommen, und bei weniger als zehn Proben oder eng benachbarten Horizonten gilt die Schätzung als unzuverlässig \cite[S.~30]{DIN50100}. Für kleine Stichproben empfiehlt die Norm, die Standardabweichung für Linien kleiner Ausfallwahrscheinlichkeit aus der Literatur zu nehmen. Extrapolationen über den Bereich zwischen $10$\,\% und $90$\,\% hinaus gelten als unsicher \cite[S.~29 und 63]{DIN50100}.

VDI 2736 Blatt 4 hält eine Prüflingsanzahl von über 30 je Versuch für eine statistische Absicherung nötig, wirtschaftlich jedoch nicht vertretbar \cite[Abschnitt~4, S.~8 bis 9]{VDI2736-4}. Mit den geforderten drei Wiederholungen je Parameterpaar hat der Kennwert den Charakter einer Schätzung \cite[S.~9]{VDI2736-4}. Die Richtlinie rechnet mit den dekadischen Logarithmen der Lastwechselzahlen und schätzt zuerst die Werte für $50$\,\%, daraus die Werte für $10$\,\% Schadenswahrscheinlichkeit \cite[S.~9]{VDI2736-4}. Vorausgesetzt werden eine Normalverteilung und ein für alle Messpunkte gleiches Verhältnis von Streuung zu Mittelwert der logarithmierten Werte \cite[S.~9]{VDI2736-4}:
\begin{equation}
    \begin{aligned}
        L_{ij}     & = \lg N_{ij} , \qquad
        L_{i,50\%} = \frac{1}{n_{i}} \sum_{j=1}^{n_{i}} L_{ij} , \qquad
        \bar{s} = \frac{1}{n_{\mathrm{p}}} \sum_{i=1}^{n_{\mathrm{p}}} \frac{s_{i}}{L_{i,50\%}} , \\
        L_{i,10\%} & = \left(1 - 1{,}28\,\bar{s}\right) L_{i,50\%} , \qquad
        N_{i,10\%} = 10^{L_{i,10\%}} ,
    \end{aligned}
    \label{eq:vdi2736_10prozent}
\end{equation}
mit der Standardabweichung $s_{i}$ der $n_{i}$ Wiederholungen des $i$-ten Parameterpaars und der Zahl $n_{\mathrm{p}}$ der Paare. Das sind die Gleichungen (10) bis (15) der Richtlinie \cite[Abschnitt~4, S.~9 bis 10]{VDI2736-4}. Die Mittelung der normierten Standardabweichungen über alle Niveaus soll die Schätzung verbessern, unterstellt aber eine Lastunabhängigkeit der Streuung \cite[S.~22]{Has18}. Zur Ausfallwahrscheinlichkeit der Festigkeitswerte in Blatt 2 macht die Richtlinie keine Aussage \cite[S.~22]{Has18}. Blatt 4 weist die aus Versuchen gewonnenen Punkte der Wöhlerlinie als Werte für $10$\,\% Schadenswahrscheinlichkeit aus \cite[S.~7]{VDI2736-4}. Ob die tabellierten Werte für PA66 und POM auf diesem Weg entstanden sind, steht in keinem der beiden Blätter. Illenberger vergleicht den Umrechnungsfaktor von $50$\,\% auf $10$\,\% nach Blatt 4 mit seiner Weibull-Auswertung und erhält $0{,}73$ gegenüber $0{,}58$ \cite[Tabelle~7-3, S.~99]{Ill22}. Für Stahlräder gelten die Festigkeitswerte nach ISO 6336 für $1$\,\% Ausfallwahrscheinlichkeit, die Hasl aus der logarithmischen Standardabweichung mit dem Quantil $2{,}33$ schätzt \cite[Abschnitt~2.1.5, S.~9]{Has18}. Weil seine Lastmomente versuchsspezifisch ausgewertet werden und keine mehrfach belegten Niveaus bilden, bezieht er die Lastspielzahl jedes Einzelversuchs auf den Funktionswert der Ausgleichsgeraden \cite[Abschnitt~8.1, S.~109 bis 110]{Has18}. Die Standardabweichung über alle bezogenen Werte einer Variante entspricht dann einer Parallelverschiebung der Geraden \cite[S.~110]{Has18}. Um bei sehr kleiner Streuung auf der sicheren Seite zu bleiben, begrenzt er die logarithmische Standardabweichung nach unten auf $0{,}1$ \cite[S.~118]{Has18}.

Als Alternative zur logarithmischen Normalverteilung dient für Ermüdungsversuche die Weibull-Verteilung. Illenberger sortiert die Versuchspunkte nach ihrer Laufzeit, ordnet dem $j$-ten von $n$ Punkten die Ausfallwahrscheinlichkeit $P_{\mathrm{A}} = j/(n+1)$ zu und trägt sie in ein Weibull-Netz ein \cite[Abschnitt~2.3.7, S.~19]{Ill22}. Die Ausgleichsgerade liefert die Lebensdauer für $50$\,\% und die Umrechnung auf kleinere Ausfallwahrscheinlichkeiten \cite[S.~19]{Ill22}. Fiederling berechnet die Ausfallwahrscheinlichkeit eines Versuchspunkts auf einem Lastniveau nach Rossow \cite[Abschnitt~2.5, S.~34]{Fie23}:
\begin{equation}
    P_{\mathrm{A}} = \frac{3\,j - 1}{3\,n + 1} ,
    \label{eq:rossow}
\end{equation}
eine Näherung, deren Ergebnisse DIN 50100 als vergleichbar mit der korrigierten Standardabweichung bezeichnet \cite[Abschnitt~10.1, S.~53]{DIN50100}. Etwa 200 Versuchspunkte zur Grübchentragfähigkeit von PEEK-Rädern ordnen sich auf den bezogenen Lastspielzahlen bis unter $2$\,\% Ausfallwahrscheinlichkeit einer Weibull-Geraden zu \cite[Abschnitt~7.8.1, S.~95 bis 96]{Ill22}. Der Formparameter beträgt $3{,}46$, die Lastspielzahl für $63{,}2$\,\% Ausfälle das $1{,}12$-Fache von $N_{50\%}$ \cite[S.~96]{Ill22}. Die Faktoren von $50$\,\% auf $10$\,\% und auf $1$\,\% betragen dort $0{,}58$ und $0{,}30$, für einsatzgehärtete Stahlräder nach der dort zitierten Auswertung von Stahl $0{,}54$ und $0{,}26$ \cite[Tabelle~7-2, S.~96]{Ill22}. Da die Weibull-Verteilung mit dem Formparameter $3{,}602$ näherungsweise einer Normalverteilung entspricht, lässt sich das Streuverhalten dieser Versuche auch mit der Normalverteilung beschreiben \cite[S.~97]{Ill22}. Für die Zahnfuß-Zeitfestigkeit von Stahlrädern gilt nach der von Fiederling wiedergegebenen Auswertung von Stahl die logarithmische Normalverteilung, deren Anwendbarkeit auf Kunststoffräder dort vermutet wird \cite[Abschnitt~2.5, S.~33]{Fie23}.

Erreicht ein Versuch die Grenzschwingspielzahl ohne Ausfall, zählt er als Durchläufer. Die Grenzschwingspielzahl ist zu dokumentieren, und eine Extrapolation der Wöhlerlinie über sie hinaus ist nicht zulässig \cite[S.~16 und 62]{DIN50100}. In die Zeitfestigkeitsgerade gehen nur Ergebnisse oberhalb des höchsten Lasthorizonts mit einem Durchläufer ein \cite[Abschnitt~10.2.1, S.~54]{DIN50100}. Eine Langzeitfestigkeit aus ein bis zwei Durchläufern oder aus der Verlängerung der Zeitfestigkeitsgeraden bis zu einer vermuteten Knickschwingspielzahl trägt nach der Norm ein erhebliches Risiko der Fehleinschätzung \cite[Abschnitt~10.3.1, S.~59]{DIN50100}. Illenberger wertet Prüfräder, die das Schadenskriterium bis $50 \cdot 10^{6}$ Lastwechsel nicht erreichen, als Durchläufer \cite[S.~66]{Ill22}.

% Herkunft:
% Absatz 1 (Kennwertermittlung): VDI2736-4.txt S. 5 (Abschnitt 3: Parameter von Blatt 2 durch Versuche an realen
%   Zahnrädern), S. 6 (feste Temperatur, Drehzahl, Drehmoment; 3.1: np Paare, eine Temperatur, vier Drehmomente,
%   drei Wiederholungen; Bruch oder Grübchenmenge; Schadensart erst nach dem Lauf), S. 7 (3.3: Gl. (4) analog
%   Blatt 2 Gl. (10); Umrechnung nach Abschnitt 4 auf Ni,10 %; Tripel als Punkt der Wöhlerlinie), S. 20 (8.4:
%   hohes Startmoment, 10^5 bis 10^6, Durchläufer verlängern).
% Absatz 2 (Prüfstände, Versagen, Temperatur): VDI2736-4.txt S. 11 (5.2 drehweich,
%   Momentmessung, Abfall durch Verschleiß oder Verformung); Hasl2018.txt Druckseite 91 (7.2: Loss-of-drive,
%   Setzeffekte, Moment vor dem Ausfall); MA-0296.txt Druckseite 65 (6.1: Totalausfall, schleichendes Risswachstum,
%   Abfall > 1 Nm über 100 000 LW, Abbildung 53); VDI2736-4.txt S. 16 (8.2 Thermografie/Thermoelement,
%   Emissionskoeffizient; 8.2.2 Öltemperatur), S. 16 und 18 (Öl erschwert optische Messung; Zuverlässigkeit
%   des Flankentemperaturkennwerts für Stirnräder nicht nachgewiesen).
% Absatz 3 (Pulsator, Korrelation): MA-0296.txt Druckseite 12 (2.2.2), 14 bis 15 (2.2.3 Unterschiede, Faktoren nur
%   für Stahl), 105 bis 106 (8.1: Linien links und steiler, 12 Nm: 5 452 210 gegen 2 058 556 LW, Faktor ≈ 2,6);
%   VDI2736-4.txt S. 12 (5.3); Hasl2018.txt Druckseite 61 bis 62 (5.1: Störgrößen, Aufwand, Laufversuch =
%   Anwendungsfall).
% Absatz 4 (Wöhlerlinie ohne Dauerfestigkeit): VDI2736-2.txt S. 14 bis 15 (3.2.2: Bild 3, Fußnote 8, SFmin 2,0 bei
%   NL = 10^8; Tabelle 5 POM 80 °C: 50/35/26/20 N/mm²); Hasl2018.txt Druckseite 107 (Darstellung), 181 (13.7.1
%   doppeltlogarithmisch, Bild 84), 176 (13.5: Frage nach Dauerfestigkeit); MA-0296.txt Druckseite 82 bis 83 (6.3.1:
%   40 Mio. LW bei 8 Nm, keine Dauerfestigkeit, These Corten/Dolan bzw. Pogačnik/Beermann bis 10^10 LW);
%   Stommel2018.txt Druckseite 328 (6.1.4 Grenzspannungsamplitude NG = 10^7), 206 bis 207 (3.2.1 Basquin Gl.
%   (3.26), (3.27), Neigungsexponent k); DIN50100.txt S. 16 (Abschnitt 5: Langzeitfestigkeit statt
%   Dauerfestigkeit, NG dokumentieren, Darstellung nur zwischen den belegten Horizonten).
% Absatz 5 (DIN 50100 Grundlagen): DIN50100.txt S. 9 (Abschnitt 1 Anwendungsbereich), S. 18 bis 19 (Abschnitt 5:
%   Verteilungsfunktionen, lg-Normalverteilung, Standardabweichung), S. 54 (10.2.1 Mittelwert mit n = 5 bis 15),
%   S. 26 (9.2.1 Perlenschnur gegen Horizonten), S. 28 (9.2.2 Regression, mehrere Versuche je Horizont),
%   S. 57 (10.2.3 mehr als zwei Horizonte).
% Absatz 6 (Gleichungen DIN 50100): Gl. (30) S. 55; Gl. (36), (37), (38) S. 56; Gl. (39), (40), (43) S. 57;
%   Gl. (67) und Tabelle 11 S. 63; parallele Linien S. 63; slgN 0,10 bis 0,30, 0,20, n < 10 S. 30; Literaturwerte und
%   Extrapolation S. 63 und S. 29. Die Streuspanne T_N (Gl. (40) und (43) S. 57) wurde beim Kürzen gestrichen.
% Absatz 7 (VDI 2736 Blatt 4 Abschnitt 4, 10-%-Frage): VDI2736-4.txt S. 8 bis 9 (über 30 Prüflinge, drei
%   Wiederholungen, Schätzung; Logarithmen; Annahmen), S. 9 bis 10 (Gl. (10) bis (15)); Hasl2018.txt Druckseite 22
%   (keine Aussage zur Ausfallwahrscheinlichkeit in Blatt 2; Gl. (24); Lastunabhängigkeit der Streuung
%   unterstellt), 9 (2.1.5 Gl. (10) mit 2,33; DIN 3990/ISO 6336 1 %), 109 bis 110 (8.1 Gl. (51) bis (53),
%   Parallelverschiebung), 118 (slog ≥ 0,1); Illenberger2022.txt Druckseite 99 (Tabelle 7-3: 0,73 gegen 0,58);
%   VDI2736-4.txt S. 7 (Ni,10 % als Punkt der Wöhlerlinie). Die BA-Aussage (BA-0486.txt Zeilen 1180 bis 1187),
%   die Werte in Blatt 2 und 3 seien für 10 % ausgewiesen, ist eine Folgerung der BA; die Richtlinie sagt das
%   nur für die nach Blatt 4 ermittelten Punkte (S. 7 und 9, Begriff dort "Schadenswahrscheinlichkeit").
% Absatz 8 (Weibull): Illenberger2022.txt Druckseite 19 (2.3.7 Schreiber, Gl. (19), Stahl fXHN), 95 bis 96
%   (7.8.1: ca. 200 Punkte, Weibullgerade bis < 2 %, k = 3,46, T = 1,12 N50 %, f10HN 0,58, f1HN 0,30, Tabelle 7-2
%   Stahl 3,20/1,20/0,54/0,26), 97 (k = 3,602 ≈ Normalverteilung nach [McC92], Bild 7-15); MA-0296.txt
%   Druckseite 33 (2.5: lg-Normalverteilung nach Stahl, Vermutung für Kunststoff), 34 (Gl. (16) Rossow);
%   DIN50100.txt S. 53 (10.1 Anmerkung: korrigierte Schätzung vergleichbar mit Rossow).
% Absatz 9 (Durchläufer): DIN50100.txt S. 16 (NG, Dokumentation), S. 62 (keine Extrapolation über NG), S. 54
%   (10.2.1), S. 59 (10.3.1 Risiko der Fehleinschätzung); Illenberger2022.txt Druckseite 66 (50 Mio. LW).
%   Gestrichen beim Kürzen: Hasl2018.txt Druckseite 186 bis 187 (13.9 Bild 91, Durchläufer im Pulsator = vor
%   Schaden abgebrochen).
% Formeln: eq:zeitfestigkeitsgerade = DIN 50100 Gl. (30) S. 55 (dort lg N = lg C − k lg La; hier N_L statt N);
%   eq:perlenschnur_streuung = DIN 50100 Gl. (36) S. 56, Gl. (38) S. 56, Gl. (39) S. 57; eq:quantil = DIN 50100
%   Gl. (67) S. 63; eq:vdi2736_10prozent = VDI 2736 Blatt 4 Gl. (10) bis (15) S. 9 bis 10 (dort ni in Gl. (11) und
%   (12), nwi in Gl. (10)); eq:rossow = MA-0296 Gl. (16) S. 34 (dort PA = (3j−1)/(3n+1) = 100 % − R).
% Bildwunsch: kap2/quelle_din50100_perlenschnur aus DIN_50100_2022-12-00_DE_3337109.pdf, Druckseite 29 (PDF-Seite
%   29), Bild 8 "Wöhlerdiagramm (schematisch), Perlenschnurverfahren"; das Diagramm mit Legende ausschneiden.
%   Ersatz, falls Bild 8 ungünstig: Bild 23 auf S. 56 (Verschieben auf einen gemeinsamen Lasthorizont) oder
%   Bild 1 auf S. 17 (Wöhlerdiagramm schematisch mit Zeitfestigkeitsgerade, Streuband, Ausfall, Durchläufer).
% Bibwunsch: keine neuen Schlüssel; verwendet werden DIN50100, VDI2736-2, VDI2736-4, Has18, Ill22, Fie23, Sto18
%   (alle in litbib.bib). Die Datei bib_kap2_6_3_8.bib im Scratchpad bleibt ohne Einträge.
% Nomenklatur (neu, nicht in Nomenklatur.tex): sigma_FlimN, N/mm², Zeitschwellfestigkeit (Nennspannung), VDI 2736
%   Blatt 4 S. 5 und Blatt 2 S. 14 (prüfen, ob 2.5.1 sie schon führt); N_L, -, Lastwechselzahl, VDI 2736 Blatt 2
%   S. 14 (dort N_L) und Hasl S. XIV; R, -, Spannungsverhältnis, MA-0296 Tabelle 1 S. 15 (in 3.7 bereits verwendet,
%   nicht in Nomenklatur.tex); L_a, (Lastgröße), Lastamplitude, DIN 50100 Tabelle 1 S. 10 bis 14; C, -, Lage der
%   Zeitfestigkeitsgeraden (Achsabschnitt für La = 1), DIN 50100 S. 15 (Zeichen kollidiert mit C des
%   Richardson-Ansatzes); k, -, Neigung der Zeitfestigkeitsgeraden, DIN 50100 S. 15 (steht als Neigungsexponent
%   der Wöhlerlinie nach Stommel bereits drin); n, -, Stichprobenumfang, DIN 50100 Tabelle 1 (kollidiert mit n
%   Drehzahl); s_lgN, -, Standardabweichung der Logarithmen der Schwingspielzahlen, DIN 50100 Tabelle 1 S. 13
%   (mit Tilde: Schätzung im Perlenschnurverfahren S. 56; Index korr S. 57); P_A, %, Ausfallwahrscheinlichkeit,
%   DIN 50100 Tabelle 1 S. 12
%   (VDI 2736 Blatt 4: Schadenswahrscheinlichkeit); u, -, Quantil für Auftretenswahrscheinlichkeit bei
%   Standardnormalverteilung, DIN 50100 Tabelle 1 S. 14; L_ij, -, log10 N_ij, VDI 2736 Blatt 4 S. 3; L_i,50 % und
%   L_i,10 %, -, logarithmierte Schätzung von N_i,50 % bzw. N_i,10 %, VDI 2736 Blatt 4 S. 3; s_i und s quer, -,
%   Standardabweichung des i-ten Versuchs und normierter Mittelwert, VDI 2736 Blatt 4 S. 9 bis 10; n_i, n_p, -,
%   Anzahl Wiederholungen bzw. Anzahl Parameterpaare, VDI 2736 Blatt 4 S. 9; j, -, Ordnungszahl, MA-0296 S. 34;
%   N_G steht bereits (Grenzlastspielzahl), hier auch Grenzschwingspielzahl nach DIN 50100 S. 12.
% Offen: Die BA-Aussage zur 10-%-Annahme (Zeilen 1180 bis 1187) ist auf den Wortlaut der Richtlinie umgestellt
%   (Absatz 7). Die Autoren Stahl [Sta99], Rossow [Ros64], Schreiber, Corten/Dolan und Pogačnik/Beermann liegen
%   nicht in 00_literatur; sie werden nur über Ill22, Fie23 und Has18 wiedergegeben. ISO 6336 (1 %) ist nach
%   Hasl S. 9 zitiert, nicht an ISO6336-1.txt geprüft. Die Roadmap nennt für den Anriss als Versagen des
%   Nachweises Stommel S. 221; der Satz steht auf S. 204 (Abschnitt 3.2), S. 221 beginnt 3.3.1.
